\clearpage
\section{EXAMPLE OF PLAY: GAME TURN 1}

This example follows the \textit{First Shots} scenario and its first-play
simplifications. Political rolls and Appeals are omitted, both players receive
5 extra RU, and hands are revealed. The named die rolls are illustrative; no
player chooses combat results in an actual game.

The players use a deliberately aggressive but legal deployment. Terra places
its three Beam Scouts at Junction and its other units at connected Terran
sites. The Imperium places one Mixed Scout and one Mixed Destroyer at
Shuruppak, with its remaining units elsewhere on its connected network. No
Base or Ground Unit is at Shuruppak. All starting ships are supplied.

Terra's revealed hand contains \textit{Tactics} and two cards that will not be
used this turn. The Imperial player has two cards but saves both.

\subsection{A. Advance the Clocks}

The Terran Turn marker advances from 0 to 1 and the Imperial marker falls from
11 to 10. Glory remains 5. Neither marker is checked for victory until the End
Phase.

\subsection{B. Joint Logistics Phase}

\subsubsection{Income}

Every starting site can trace through friendly Worlds and Outposts to its
capital. Terra receives 8 RU for each of three connected Worlds and 1 RU for
each of three connected Outposts. The Imperium receives its 10 RU provincial
budget and 3 RU for each of four connected Worlds; Imperial Outposts produce
no RU. Each side then adds the learning-game bonus.

\begin{tabularx}{\columnwidth}{@{}Xrr@{}}
\toprule
Turn 1 income & Terra & Imperium \\
\midrule
Worlds & $3\times8=24$ & $4\times3=12$ \\
Outposts & $3\times1=3$ & 0 \\
Provincial budget & 0 & 10 \\
Learning-game bonus & 5 & 5 \\
\midrule
RU before maintenance & 32 & 27 \\
\bottomrule
\end{tabularx}

\subsubsection{Maintenance}

Maintenance equals Hull Size. Bases, Ground Units, and Fleet Trains are free.
All ships are supplied, so there is no surcharge.

\begin{tabularx}{\columnwidth}{@{}Xrr@{}}
\toprule
Starting ships & Terra & Imperium \\
\midrule
SC, DD & $3+2$ & $2+3$ \\
CA, BC & $2+0$ & $4+2$ \\
CV, Fighters & $1+2$ & $1+2$ \\
Transports & 2 & 1 \\
\midrule
Maintenance paid & 12 & 15 \\
RU remaining & 20 & 12 \\
\bottomrule
\end{tabularx}

\subsubsection{Cards, Arrivals, and Production}

Both hands already equal their minimum size, so neither side draws or buys a
card. Nothing is scheduled to arrive on Game Turn 1.

Terra buys one Destroyer for its printed cost of 9 RU and assigns it a
\textbf{Beam Package}. Terra also buys a 5-RU Fleet Train and a 4-RU Outpost,
saving 2 RU. The Imperium buys a 9-RU Destroyer, assigns it a
\textbf{Mixed Package}, and saves 3 RU. All four purchases are placed on Game
Turn 2; none can move or fight this turn.

\begin{center}
\begin{tikzpicture}[
  box/.style={draw,rounded corners=2pt,align=center,minimum width=2.2cm,
    minimum height=0.75cm,font=\scriptsize},
  arr/.style={-{Latex[length=2mm]},thick}
]
  \node[box,fill=Cyan!12] (ti) {Terra\\32 RU};
  \node[box,fill=Cyan!12,right=0.35cm of ti] (tm) {$-12$ maintenance\\20 RU};
  \node[box,fill=Cyan!12,below=0.3cm of tm] (tb) {$-9$ DD, $-5$ Train\\$-4$ Outpost; save 2};
  \draw[arr,Cyan!70!black] (ti) -- (tm);
  \draw[arr,Cyan!70!black] (tm) -- (tb);

  \node[box,fill=BrickRed!12,below=0.3cm of ti] (ii) {Imperium\\27 RU};
  \node[box,fill=BrickRed!12,below=0.3cm of ii] (im) {$-15$ maintenance\\12 RU};
  \node[box,fill=BrickRed!12,right=0.35cm of im] (ib) {$-9$ DD\\save 3};
  \draw[arr,BrickRed!75!black] (ii) -- (im);
  \draw[arr,BrickRed!75!black] (im) -- (ib);
\end{tikzpicture}
\end{center}

\subsection{C. Terran Player Turn}

\subsubsection{First Movement: the Shuruppak Raid}

The three Beam Scouts at Junction form one stack and jump through Southern
Cross to Sirius. There they split: one Scout stops at Sirius as a reserve while
two continue through Mackhashi and Tau Ceti before entering Shuruppak on their
fifth jump. Each intermediate system is empty. This movement is legal because
a stack composed entirely of Scouts may refuel in any system. An ordinary
mixed fleet would have to stop one jump beyond its last friendly refuelling
site.

Entering Shuruppak reveals the Imperial Mixed Scout and Mixed Destroyer and
ends the two-Scout raiding stack's movement. Space combat is mandatory.

\begin{center}
\begin{tikzpicture}[
  sys/.style={circle,draw,minimum size=6.5mm,inner sep=0pt,font=\scriptsize},
  route/.style={draw=ForestGreen!75!black,line width=1.1pt,
    -{Latex[length=1.7mm]}},
  lab/.style={font=\scriptsize,align=center}
]
  \node[sys,fill=Cyan!18] (j) at (0,0) {0};
  \node[sys] (sc) at (1.15,0.55) {1};
  \node[sys] (si) at (2.3,0) {2};
  \node[sys] (ma) at (3.45,0.55) {3};
  \node[sys] (tc) at (4.6,0) {4};
  \node[sys,fill=BrickRed!18] (sh) at (5.75,0.55) {5};
  \draw[route] (j)--(sc);
  \draw[route] (sc)--(si);
  \draw[route] (si)--(ma);
  \draw[route] (ma)--(tc);
  \draw[route] (tc)--(sh);
  \node[lab,below=1mm of j] {Junction\\3 Beam SC};
  \node[lab,above=1mm of sc] {Southern\\Cross};
  \node[lab,below=1mm of si] {Sirius\\1 SC stops};
  \node[lab,above=1mm of ma] {Mackhashi};
  \node[lab,below=1mm of tc] {Tau Ceti};
  \node[lab,above=1mm of sh] {Shuruppak\\2 SC vs SC+DD};
\end{tikzpicture}
\end{center}

The numbered circles show jumps, not movement costs. The route drawing is a
schematic of the printed map rather than a replacement for it.

\subsubsection{First Combat: Round 1 at Long Range}

Terra is the attacker. The initial range is Long. Each side has two ships, so
there is no screening or Fleet Size Bonus.

Before the first firing round, Terra plays \textit{Tactics}, raising its
Tactics from 0 to 1 for this entire battle. Its once-per-round attack reroll is
available, although Terra will not need it in the example.

One Terran Beam Scout performs Point Defense for its group. It forfeits
offensive fire, which costs nothing because Beam ships cannot fire at Long
Range, and the group gains +1 Defense against missiles.

The Imperial Mixed Destroyer has printed Attack 4, reduced to 3 for Mixed
missile fire at Long Range. Against Scout Defense 0 and Point Defense +1, its
Hit Number is $3-0-1=2$; a roll of 4 misses. The Mixed Scout's adjusted Attack
is 2, giving a Hit Number of 1. It rolls a natural 1 and destroys one Hull-1
Terran Scout. Terra removes the non-escort Scout. Its remaining Beam ship has
no legal fire, so the round ends with one Terran ship against two Imperial
ships.

\subsubsection{Round 2: Terra Closes}

Both players select a Scout Picket, worth +2 Maneuver. Terra rolls 7, for
$7+2+1+1=11$; the final +1 is for having the smaller non-Fighter fleet. The
Imperium rolls 4, for $4+2+1=7$. Terra chooses Close Range.

The Imperium now has two ships to one. It screens its Destroyer, forcing the
Terran Scout to target the unscreened Scout. Screening leaves only one
unscreened Imperial ship, however, so the Imperium gives up what would
otherwise have been a $2{:}1$ Fleet Size Bonus.

At Close Range beams fire by class before short-range missiles. The Mixed
Destroyer fires at Class D with adjusted Attack 3 but rolls 6 and misses. At
Class E both sides have Tactics 1, so the Imperial defender fires first. Its
Mixed Scout has adjusted Attack 2 and rolls 5, missing.

The Terran Beam Scout attacks at its full printed Attack 3. The Destroyer is
screened, so it must target the Defense-0 Imperial Scout. It rolls 2 and
destroys that ship with one hit. Because the Imperial Scout chose beam fire,
it could not also make a short-range missile attack had it survived.

\subsubsection{Round 3: Missile Range Restored}

Terra rolls $3+2+1=6$ for Range Control. The surviving Imperial Destroyer uses
its +1 Picket rating and Tactics 1. It rolls $8+1+1=10$ and chooses Long Range.

Neither side can screen or claim Fleet Size. The Terran Scout again performs
Point Defense. The Imperial Destroyer declares High-Intensity Missile Fire:
printed Attack 4, $-1$ for its Mixed Package, and $+1$ for high intensity.
After Point Defense, its Hit Number is $4-0-1=3$. It rolls 2 and destroys the
last raiding Scout before that Scout's Class E opportunity to retreat. The
Imperium retains Shuruppak.

\begin{center}
\begin{tikzpicture}[
  rnd/.style={draw,rounded corners=2pt,align=left,text width=5.8cm,
    inner sep=3pt,font=\scriptsize},
  flow/.style={-{Latex[length=1.7mm]},thick}
]
  \node[rnd,fill=BrickRed!10] (r1) {\textbf{Round 1 --- Long:}
    Imperial missiles destroy 1 of 2 raiding Beam Scouts.};
  \node[rnd,fill=Cyan!10,below=2mm of r1] (r2) {\textbf{Round 2 --- Close:}
    Terra wins Range Control and destroys the unscreened Mixed Scout.};
  \node[rnd,fill=BrickRed!10,below=2mm of r2] (r3) {\textbf{Round 3 --- Long:}
    the Mixed DD wins range and destroys the last raider with HIMF.};
  \draw[flow] (r1)--(r2);
  \draw[flow] (r2)--(r3);
\end{tikzpicture}
\end{center}

\subsubsection{Reaction Movement and Combat}

The Imperial player uses the surviving Mixed Destroyer as the one reacting
stack. It jumps from Shuruppak through Tau Ceti and Mackhashi to Sirius, where
it must stop upon encountering the reserve Terran Scout. Three jumps exactly
equal the Imperial Reaction Allowance. The reacting Destroyer is the only
Imperial unit permitted to attack in the Reaction Combat Phase.
Its Missile Depleted marker was removed at the end of the preceding Combat
Phase, so it may make another missile attack in this new phase.

The new battle begins at Long Range. The Terran Scout uses Point Defense; the
Mixed Destroyer's Attack 3 missile therefore has a Hit Number of 2 and its roll
of 7 misses. In round two, Terra rolls $9+2=11$ for Range Control; the Imperium
rolls $4+1+1=6$. Terra chooses Close. The Class D Mixed Destroyer fires first
with adjusted Attack 3, but rolls 8 and misses. At Class E the Terran Beam
Scout rolls a natural 1 and destroys the Hull-1 Destroyer. The attempted
pursuit has cost the Imperium its remaining forward ship.

\subsubsection{Second Movement and Combat}

The surviving reserve Scout returns from Sirius through Southern Cross to
Junction in the Second Movement Phase. Terra has no contested system in its
Second Combat Phase.

The Shuruppak Outpost remains Imperial because no Terran ship remains in orbit.
Even if a Scout had survived, bombardment alone could only remove temporary
Militia; it could not capture the surface. Capturing the Outpost requires a
later fleet with a Transport and surviving troops to pay the occupation cost.

\subsection{D. Imperial Player Turn}

The Imperium uses its two Movement Phases to shift its Carrier, Fighters, and
main surface force toward the Shuruppak corridor. A Fleet Train can move only
one jump in each phase and provides refuelling only when it has not moved in
that particular phase, so the main fleet does not reach Terran space this
turn. Terra declines Reaction Movement. No further combat occurs.

This illustrates why the Scout raid was possible while a combined fleet was
not: self-refuelling is a Scout ability, not a property of every stack that
contains a Scout.

\subsection{E. End Phase}

Temporary combat markers are removed. No World was neutralized and no site
changed ownership, so Glory remains 5. The Terran marker is 1 and the Imperial
marker is 10; neither has reached Glory. Game Turn 2 will begin with the new
Destroyers, Terran Fleet Train, and Terran unplaced Outpost arriving during
Completed Production.

\subsection{Turn 1 Lessons}

\begin{itemize}
  \item Income is collected before maintenance and new production; newly
    purchased units do not appear immediately.
  \item Armament Packages create different preferred ranges without replacing
    the printed \SpaceEmpires\ combat values.
  \item Point Defense can be valuable even when its escort had no legal
    offensive shot.
  \item Range Control can reverse the advantage from round to round.
  \item Reaction Movement creates a real counterstroke, but only reacting
    ships may attack in its Combat Phase.
  \item Space victory and surface conquest are separate achievements.
\end{itemize}
